\chapter{Перенесення частинок, енергії та домішок}\label{ch:transport}

Рівновага (розд.~\ref{ch:equilibrium}) задає магнітні поверхні, а перенесення (transport) визначає, як на них повільно змінюються профілі $n_s$, $T_s$ і струму. Поперечні потоки мають три джерела: класичне (кулонівські зіткнення, крок~--- ларморів радіус), неокласичне (тороїдальна геометрія, крок~--- ширина банана) та аномальне (турбулентність). У ядрі переважає аномальне. Потоки, $\Er$ і перенесення імпульсу розглянуто в розд.~\ref{ch:flows}.

\section{Кулонівські зіткнення і неокласичні режими}\label{sec:05-coll}

У корпусі частоти зіткнень у СІ задає час обміну енергією між сортами $s$ і $s'$~\cite[с.~3, (14), (20)]{Seto2014}:
\begin{equation}\label{eq:05-tau}
  \tau_{ss'}=\frac{3\sqrt2\,\pi^{3/2}\varepsilon_0^2\,m_sm_{s'}}{n_{s'}e^4Z_s^2Z_{s'}^2\lnL}\left(\frac{T_s}{m_s}+\frac{T_{s'}}{m_{s'}}\right)^{3/2},\qquad
  P_{\Delta s}=\sum_{s'}\frac32\,n_s\frac{T_{s'}-T_s}{\tau_{ss'}} .
\end{equation}
Для пари $e$--$i$ ($T_e/m_e\gg T_i/m_i$) маємо $\tau_{ei}=(m_i/2m_e)\,\tau_e$, де брагінський час електрон-іонних зіткнень
$\tau_e=6\sqrt2\pi^{3/2}\varepsilon_0^2m_e^{1/2}T_e^{3/2}/(\lnL\,e^4Z^2n_i)$. Тоді $P_{\Delta e}=-3(m_e/m_i)\,n_e(T_e-T_i)/\tau_e$; саме в такому вигляді обмін енергією входить у баланс Тищенка \eqref{eq:05-tysh}~\cite[с.~109, (4.54)]{Tyshchenko2021}. Швидкі іони (пучок, α) гальмуються на електронах за час $\tau_{sb}=6\pi\sqrt{2\pi}\varepsilon_0^2m_bT_e^{3/2}/(n_eZ_b^2m_e^{1/2}e^4\lnL)$~\cite[с.~4, (8)]{Fukuyama1992} (розд.~\ref{ch:heating}).

\begin{outofcorpus}[$\lnL$, зіткненнєвість $\nustar$ і три режими]
Кулонівський логарифм $\lnL\approx31{,}3-\ln(\sqrt{n_e}/T_e)$ ($n_e$ у м$^{-3}$, $T_e$ в еВ): 17--18 у ядрі, 13--14 на краю. Іонний час $\tau_i=12\pi^{3/2}\varepsilon_0^2m_i^{1/2}T_i^{3/2}/(\lnL\,e^4Z^4n_i)$ довший за $\tau_e$ приблизно в $\sqrt{m_i/m_e}$ разів.

\emph{Захоплена частинка} (п.~\ref{sec:02-banana}) займає конус пітч-кутів шириною $\Delta\xi\sim\sqrt\varepsilon$. Розсіяння на кут $\sqrt\varepsilon$ відбувається за час $\varepsilon/\nu$, тож ефективна частота $\nu_{\mathrm{eff}}=\nu/\varepsilon$. Баунс-частота $\omega_b\sim\sqrt\varepsilon\,v_T/(qR)$, $v_T=\sqrt{T/m}$. Їхнє відношення~--- \emph{зіткненнєвість}
\begin{equation}\label{eq:05-nustar}
  \nustar=\frac{\nu_{\mathrm{eff}}}{\omega_b}=\frac{\nu\,qR}{\varepsilon^{3/2}v_T}.
\end{equation}
Режими: \textbf{банановий} $\nustar<1$ (банан встигає замкнутися); \textbf{плато} $1<\nustar<\varepsilon^{-3/2}$; \textbf{Пфірша--Шлютера} (PS) $\nustar>\varepsilon^{-3/2}$, тобто довжина вільного пробігу $v_T/\nu$ коротша за довжину зв'язку $qR$.

\emph{Оцінки випадкового блукання} $D\sim(\Delta r)^2/\Delta t$, $\rho=v_T/\Omc{s}$:
класичне $D_{\mathrm{cl}}\sim\rho^2\nu$; PS $D_{\mathrm{PS}}\approx(1+2q^2)D_{\mathrm{cl}}$ (зміщення $q\rho$ від вертикального дрейфу за час проходу довжини зв'язку); банан~--- частка $f_t\sim\sqrt\varepsilon$ частинок робить кроки $w_b\sim q\rho/\sqrt\varepsilon$ \eqref{eq:02-wb} з частотою $\nu/\varepsilon$:
\begin{equation}\label{eq:05-Dban}
  D_{\mathrm b}\sim\sqrt\varepsilon\cdot\frac{q^2\rho^2}{\varepsilon}\cdot\frac{\nu}{\varepsilon}=\frac{q^2\rho^2\nu}{\varepsilon^{3/2}},\qquad
  D_{\mathrm p}\sim\frac{q\rho^2v_T}{R}\ \text{(плато, не залежить від $\nu$)}.
\end{equation}
При $\nustar=1$ маємо $D_{\mathrm b}=D_{\mathrm p}$, а при $\nustar=\varepsilon^{-3/2}$ маємо $D_{\mathrm p}=D_{\mathrm{PS}}$: криві зшиваються неперервно. Іонна теплопровідність у банановому режимі (Хінтон--Хейзелтайн) $\chi_{i}^{\mathrm{neo}}\approx0{,}66\,q^2\rho_i^2\nu_{ii}\varepsilon^{-3/2}$; електронна менша в $\sim\sqrt{m_e/m_i}$ разів. Тому неокласика пояснює хіба що $\chi_{i}$ поблизу осі.

\emph{Приклад (JET, табл.~4.1 з~\cite[с.~86]{Tyshchenko2021}).} Ядро: $n_e=\SI{3,8e19}{м^{-3}}$, $T_e=\SI{9,9}{кеВ}$, $q=1$, $R=\SI{3}{м}$, $\varepsilon=0{,}1$ дають $\lnL=18$, $\tau_e=\SI{5,0e-4}{с}$, $\nu_{*e}\approx5\cdot10^{-3}$~--- глибоко банановий режим. Край ($n_e=\SI{2e19}{м^{-3}}$, $T_e=\SI{100}{еВ}$, $q=3{,}5$, $\varepsilon=0{,}35$): $\nu_{*e}\approx10>\varepsilon^{-3/2}=4{,}8$~--- PS.
\end{outofcorpus}

Отже, аналітичні моделі корпусу, записані для режиму PS (обертання й домішки~\cite{Sigmar1987,Wong1987}, крайовий потік~\cite{Singh2004}), стосуються краю плазми і важких домішок. У ядрі потрібна бананова неокласика, якої в корпусі немає.

\begin{outofcorpus}[бутстреп-струм і неокласичний опір]
Нехай $\dd n/\dd r<0$. У точці $r$ захоплені частинки, що йдуть «вперед», мають центри бананів на $\pm w_b/2$ від частинок, що йдуть «назад». Тому їхні густини відрізняються на $w_b\,\dd n/\dd r$, і виникає \emph{банановий струм} $j_{\mathrm b}\sim e\,(f_t w_b\,\dd n/\dd r)\,v_\parallel$ з $v_\parallel\sim\sqrt\varepsilon\,v_T$, $w_b\sim qv_T/(\sqrt\varepsilon\,\Omc{s})$. Врахувавши $\Bpol\approx\varepsilon B/q$, отримуємо $j_{\mathrm b}\sim\varepsilon^{3/2}(T/\Bpol)\,\dd n/\dd r$. Пролітні частинки отримують від захоплених імпульс із частотою $\nu/\varepsilon$, а втрачають його з частотою $\nu$. Тому в стаціонарі їхній струм у $1/\varepsilon$ разів більший:
\begin{equation}\label{eq:05-jbs}
  J_{\mathrm{BS}}\sim-\frac{\sqrt\varepsilon}{\Bpol}\frac{\dd p}{\dd r};\qquad
  J_{\mathrm{BS}}=-\frac{\sqrt\varepsilon}{\Bpol}\Big[2{,}44(T_e+T_i)\frac{\dd n}{\dd r}+0{,}69\,n\frac{\dd T_e}{\dd r}-0{,}42\,n\frac{\dd T_i}{\dd r}\Big]
\end{equation}
(друга рівність~--- стандартний результат для $\varepsilon\ll1$, $\nustar\ll1$, $Z=1$). Струм тече без зовнішнього поля. Інтегрування по перерізу дає частку $I_{\mathrm{BS}}/\Ip\sim\sqrt\varepsilon\,\betap$. Захоплені електрони не несуть струму, тому провідність зменшується: грубо $\eta_{\mathrm{neo}}\approx\eta_{\mathrm{Sp}}/(1-\sqrt\varepsilon)^2$.
\end{outofcorpus}

\section{Неокласичне замикання моментних рівнянь}\label{sec:05-closure}

Рівняння моментів (неперервність, імпульс, внутрішня енергія, повний тепловий потік) для кожного сорту~\cite[с.~3, (8)--(11)]{Seto2014} незамкнені: потрібні тертя $\vect{F}_s$, «теплове тертя» $\vect{H}_s$, тензор в'язкості $\visc_s$ і тепловий тензор в'язкості $\vect{\Theta}_s$. У моментному підході (Хіршман--Сігмар, moment approach) тертя лінійне за потоками~\cite[с.~4, (22)--(23)]{Seto2014}:
\begin{equation}\label{eq:05-fric}
  \vect{F}_s=\sum_{s'}\big(l_{11}^{ss'}\uvec_{s'}-l_{12}^{ss'}\vect{h}_{s'}\big),\qquad
  \vect{H}_s=\sum_{s'}\big(-l_{21}^{ss'}\uvec_{s'}+l_{22}^{ss'}\vect{h}_{s'}\big),\qquad \vect{h}_s\equiv\frac{2\vect{Q}^{\mathrm c}_s}{5p_s},
\end{equation}
де $\vect{Q}^{\mathrm c}_s$~--- кондуктивний тепловий потік, а $l^{ss'}_{ij}$ виражаються через матричні елементи оператора зіткнень Брагінського. У найнижчому порядку дрейфового впорядкування $\delta=\rho_i/L_\perp$ в'язкості мають форму Чу--Голдбергера--Лоу (CGL)~\cite[с.~4, (24)--(28)]{Seto2014}:
\begin{gather}
  \visc_s=\pi_{\parallel s}\Big(\unitb\unitb-\tfrac13\tens{I}\Big)+\order{\delta^2},\qquad
  \vect{\Theta}_s=\Theta_{\parallel s}\Big(\unitb\unitb-\tfrac13\tens{I}\Big)+\order{\delta^2},\label{eq:05-cgl}\\
  \begin{pmatrix}\pi_{\parallel s}\\ \Theta_{\parallel s}\end{pmatrix}=-\frac32
  \begin{pmatrix}\mu_{s1}&\mu_{s2}\\ \mu_{s2}&\mu_{s3}\end{pmatrix}
  \begin{pmatrix}W^{u}_{s}\\ W^{h}_{s}\end{pmatrix},\qquad
  W^{u}_s=2\big(\nabla_\parallel u_{\parallel s}-\uvec_s\cdot\vect{\kappa}\big),\quad
  W^{h}_s=2\big(\nabla_\parallel h_{\parallel s}-\vect{h}_s\cdot\vect{\kappa}\big),\label{eq:05-mu}
\end{gather}
$\vect{\kappa}=\unitb\cdot\grad\unitb$ (прийнято $\grad\cdot\uvec_s=0$, $\grad\cdot\vect{h}_s=0$). Ключова властивість~\eqref{eq:05-mu}: \emph{усередині LCFS} вона еквівалентна в'язкості Хіршмана в сенсі усереднених сил $\fsa{\Bvec\cdot\grad\cdot\visc_s}$, а \emph{зовні}~--- в'язкості Брагінського. Тож одна система описує ядро і SOL без зшивання~\cite[с.~4]{Seto2014}. Коефіцієнти $\mu_{sj}$ з баунс-усередненого дрейфово-кінетичного рівняння є функціями потоку (без полоїдальної залежності); на зіткненнєвій периферії \eqref{eq:05-mu} зводиться до форми Брагінського і полоїдальна залежність відновлюється.

Після підстановки рівноважних потоків $\uvec_s=\omega_{us}R^2\grad\tor+L_{us}\Bvec$, $\vect{h}_s=\omega_{hs}R^2\grad\tor+L_{hs}\Bvec$ ($\omega$~--- кутові частоти, $L$~--- пов'язані з полоїдальним потоком) і усереднення паралельних балансів імпульсу й теплового потоку отримуємо систему $2\times2$ на сорт~\cite[с.~7, (73)--(78)]{Seto2014}: в'язкість $\fsa{3(\nabla_\parallel B)^2}\,\mu_{sj}$, що діє на $(L_{us},L_{hs})$, урівноважують тертя $l^{ss'}_{ij}$ з усіма сортами \eqref{eq:05-fric} і $e_sn_s\fsa{\Epar B}$. У матричному вигляді її записано в розд.~\ref{ch:flows}, \eqref{eq:06-seto78}. Це стандартна неокласика Хіршмана--Сігмара. Для іонів вона дає полоїдальне обертання (розд.~\ref{ch:flows}). У стандартній теорії (поза корпусом) та сама система для електронів дає закон Ома $\fsa{J_\parallel B}=\sigma_{\mathrm{neo}}\fsa{\Epar B}+\fsa{J_{\mathrm{BS}}B}$, тобто \eqref{eq:05-jbs}. Зміст $\mu_{sj}$ (стандартно): тертя пролітних частинок об захоплені («магнітне накачування»); у банановому режимі $\mu\propto f_t\nu$, у PS $\mu\propto p/\nu$. У~\cite{Seto2014} $\Bvec=\grad\tor\times\grad\psi+I\grad\tor$, тобто знак $\psi$ протилежний до EFIT (див. розд.~\ref{ch:geometry}).

\section{1.5D-система перенесення}\label{sec:05-15d}

Код TASK/TR~\cite[с.~3--4, (1)--(6)]{Fukuyama1992} розв'язує рівняння для кожного сорту $s=e,D,T,\mathrm{He}$ (радіус $r$, квазінейтральність $n_e=\sum Z_in_i$):
\begin{gather}
  \frac{\pa n_s}{\pa t}=-\frac1r\frac{\pa}{\pa r}(r\Gamma_s)+S_s,\qquad \Gamma_s=v_sn_s-D_s\frac{\pa n_s}{\pa r},\label{eq:05-cont}\\
  \frac{\pa}{\pa t}\Big(\frac32n_sT_s\Big)=-\frac1r\frac{\pa}{\pa r}(rQ_s)+P_s,\qquad Q_s=v_sn_sT_s-n_s\chi_{s}\frac{\pa T_s}{\pa r}+\frac32T_s\Gamma_s,\label{eq:05-energy}\\
  \frac{\pa\Bpol}{\pa t}=\frac{\pa E_\tor}{\pa r},\qquad E_\tor=\eta\Big(\frac{1}{\mu_0r}\frac{\pa(r\Bpol)}{\pa r}-J_{\mathrm{NB}}-J_{\mathrm{BS}}\Big).\label{eq:05-Bp}
\end{gather}
Тут $\eta$~--- неокласичний опір, а $J_{\mathrm{BS}}$~--- за Хінтоном і Хейзелтайном~\cite[с.~4, 6]{Fukuyama1992}.

\textbf{Про множник $\tfrac32$ у $Q_s$.} Формулу \eqref{eq:05-energy} відтворено з джерела дослівно. Стандартно пишуть $Q_s=-n_s\chi_{s}\pa_rT_s+\tfrac52T_s\Gamma_s$ ($\tfrac52T$~--- ентальпія на частинку) з роботою $\uvec_s\cdot\grad p_s$ у правій частині; так і в~\cite[с.~3, (10), (13)]{Seto2014}: $\vect{Q}_s=\vect{Q}^{\mathrm c}_s+\tfrac52p_s\uvec_s+\dots$ Множник $\tfrac32$ еквівалентний, якщо в $P_s$ стоїть $-p_s\grad\cdot\uvec_s$. Доданок $v_sn_sT_s$ у~\cite{Fukuyama1992} за тієї самої $v_s$, що й у $\Gamma_s$, дає пінчевій конвекції множник $\tfrac52$, а дифузійній~--- $\tfrac32$; джерело цього не пояснює. У власному коді беріть $\tfrac52$ з явною роботою тиску.

Рівняння \eqref{eq:05-Bp}~--- дифузія полоїдального поля з коефіцієнтом $\eta/\mu_0\approx\SI{e-3}{м^2/с}$ при \SI{10}{кеВ}. Тому в ITER-класі ($a=\SI{2}{м}$) профіль струму релаксує повільно (джерело: «сотні секунд»; наша оцінка $\mu_0a^2/\eta\approx\SI{4e3}{с}$), і зустрічне поле після ввімкнення NBI довго зберігається~\cite[с.~10]{Fukuyama1992}. $D_s$ стале, пінч $v_s$ підібрано під профіль густини, $\chi_{s}$ дає модель дрейфових хвиль (п.~\ref{sec:05-anom}); при $\Vloop=\SI{0,06}{В}$ у запаленому стані горіння триває $\approx\SI{1600}{с}$~\cite[с.~9]{Fukuyama1992}.

\textbf{Сучасні 1.5D-коди.} В ASTRA (тренажер Fenix) теплопровідності мають вигляд~\cite[с.~4]{Fable2022}
\begin{equation}\label{eq:05-astra}
  \chi_{e,i}=\chi_{\mathrm{gB}}\,f_{e,i}(\text{геометрія},T_e/T_i,q,s,\nu,\beta,\dots)+\chi_{e,i}^{\mathrm{neo}},\qquad \chi_{\mathrm{gB}}\propto\frac{T^{3/2}\sqrt{M}}{B^2R}.
\end{equation}
Вільні коефіцієнти підганяють під базу розрядів; для частинок додають конвекцію. Струми котушок AUG збігаються з експериментом лише за \emph{зменшеного} крайового бутстреп-струму~\cite[с.~6--7]{Fable2022}. TSC для NSTX \#100920 відтворює магнітні сигнали з множником 1,5 до «стандартних» коефіцієнтів за $\Zeff=2{,}7$. «Майже ідентичне» узгодження дає множник 1,0 за $\Zeff=5$~\cite[с.~4]{Jardin2000}. Більший $\Zeff$ підвищує опір і омічний нагрів, що компенсує менший $\chi_{s}$: без незалежного виміру $\Zeff$ коефіцієнти не визначити.

\begin{corpusexample}[баланс $T_e$, $T_i$ для JET DTE1]
Тищенко~\cite[с.~108--109, (4.52)--(4.54)]{Tyshchenko2021} розв'язує
\begin{equation}\label{eq:05-tysh}
  \frac32n_s\frac{\pa T_s}{\pa t}-\frac1r\frac{\pa}{\pa r}\Big(rn_s\chi_{s}\frac{\pa T_s}{\pa r}\Big)=\pm P_{ei}+P^{s}_{\mathrm{nbi}}+P^{s}_{\alpha}+\begin{cases}-P_{\mathrm{br}}, & s=e,\\ P_{\mathrm w}, & s=i,\end{cases}
\end{equation}
($+$ для електронів) з обміном $P_{ei}=P_{\Delta e}$ за \eqref{eq:05-tau}, $P_{\mathrm{br}}\propto n_en_i\sqrt{T_e}$, профілями NBI з TRANSP і сталими $\chi_{i}=\SI{5000}{см^2/с}=\SI{0,50}{м^2/с}$, $\chi_{e}=\SI{2700}{см^2/с}=\SI{0,27}{м^2/с}$. Їх підібрано під D-розряд \#41069, 10~МВт NBI~\cite[с.~111]{Tyshchenko2021} (там для $\chi_{e}$ описка «см$^2$\,м$^{-1}$»; у підписах с.~113--114~--- см$^2$\,с$^{-1}$). Доданок $P_{\mathrm w}$~--- частка $\fSC$ потужності α, передана хвилями (каналювання, розд.~\ref{ch:heating}).

\begin{center}\small\setlength{\tabcolsep}{4pt}
\begin{tabular}{lccc|lcc}
\toprule
Модель~\cite[с.~114, табл.~4.2]{Tyshchenko2021} & $T_{i0}$, кеВ & $T_{e0}$, кеВ & $T_{i0}-T_{e0}$ & Експеримент~\cite[с.~86]{Tyshchenko2021} & $T_{i0}$ & $T_{e0}$\\
\midrule
NBI & 13,35 & 9,90 & 3,45 & \#41069 (D) & 13,3 & 9,9\\
NBI+α & 15,71 & 12,96 & 2,75 & \#42856 (50\,\% T) & 16,5 & 11,5\\
NBI+α+SC, $\fSC=0{,}2$ & 16,01 & 12,53 & 3,48 & \#42847 (75\,\% T) & 17,5 & 12,0\\
NBI+α+SC, $\fSC=0{,}3$ & 16,15 & 12,32 & 3,83 & & & \\
\bottomrule
\end{tabular}
\end{center}
З тими самими $\chi_{s}$ α-нагрів перегріває електрони, а каналювання збільшує $T_i-T_e$. Проте навіть $\fSC=0{,}3$ не дає виміряних $T_i>\SI{16,2}{кеВ}$, $T_e\lesssim\SI{12}{кеВ}$; більшу частку автор відкидає, бо розрахунки для FMM дають лише 15--20\,\% енергії α. Структуру FMM на JET не виміряно, тож лишаються й інші пояснення (стабілізація ITG, ізотопний ефект, нестаціонарність)~\cite[с.~115, 117]{Tyshchenko2021}.
\end{corpusexample}

\section{Аномальне перенесення}\label{sec:05-anom}

Виміряна електронна температуропровідність на два порядки перевищує неокласичну. Іонна~--- лише в кілька разів, але має іншу залежність від параметрів~\cite[с.~5]{Fukuyama1992}. У TASK/TR коефіцієнти задає турбулентність дрейфових хвиль за теорією довжини перемішування (mixing length): $\chi_{s}\propto\gamma/k_\perp^2$, де $\gamma$~--- лінійний інкремент, $k_\perp$~--- поперечне хвильове число. Внески дають моди захоплених і пролітних електронів (беззіткненнєві й дисипативні) та ITG-мода (ion temperature gradient; формули Домінгеса--Волтца). Множник порядку одиниці нормують на скейлінг утримання (п.~\ref{sec:05-scal})~\cite[с.~5, 8]{Fukuyama1992}. У~\cite[с.~3, (21)]{Seto2014} турбулентність входить у рідинні рівняння як квазілінійна сила $\vect{F}^{\mathrm{QL}}_s$. Вона пропорційна квадрату амплітуди хвиль і автоматично дає амбіполярний потік частинок.

\begin{outofcorpus}[від довжини перемішування до гіро-Бома]
Вихор розміром $k_\perp^{-1}$ перевертається за час $\gamma^{-1}$, тож $D\sim\gamma/k_\perp^2$. Насичення настає, коли градієнт збурення зрівнюється з рівноважним: $\tilde n/n\sim1/(k_\perp L)$, $L$~--- масштаб градієнта. Для ITG $k_\perp\rho_i\sim1$, $\gamma\sim v_{Ti}/L$, тож
\begin{equation}\label{eq:05-gB}
  \chi_{\mathrm{gB}}\sim\frac{\rho_i^2v_{Ti}}{L}=\frac{\rho_i}{L}\,\chi_{\mathrm B},\qquad \chi_{\mathrm B}\sim\rho_iv_{Ti}=\frac{T}{eB};\qquad
  L\sim R:\ \chi_{\mathrm{gB}}\sim\frac{\sqrt{m_i}\,T^{3/2}}{e^2B^2R},
\end{equation}
що збігається з \eqref{eq:05-astra}. Для JET ($T=\SI{5}{кеВ}$, D, $B=\SI{3,45}{Тл}$, $R=\SI{3}{м}$) $\chi_{\mathrm{gB}}\approx\SI{1,4}{м^2/с}$. Це той самий порядок, що й $\chi_{i,e}$ Тищенка; сам $\chi_i=\SI{0,5}{м^2/с}$ у $\sim70$ разів більший за банановий $\chi_i^{\mathrm{neo}}$ (задача~1). ITG вмикається, коли $R/L_{T_i}$ перевищує критичне значення $\sim4$--6. Вище порогу потік різко зростає («жорсткість» профілів), тому функції $f_{e,i}$ в \eqref{eq:05-astra} залежать від $T_e/T_i$, ширу $s$ і $q$. Бомівський $\chi_{\mathrm B}$ дає $\tauE\Omc{i}\propto\rho_*^{-2}$, гіро-бомівський~--- $\rho_*^{-3}$ ($\rho_*=\rho_i/a$).
\end{outofcorpus}

\begin{corpusexample}[глобальна гірокінетика GTC: розтікання турбулентності]
Глобальний δf-код частинок у комірках GTC у загальній геометрії DIII-D~\cite[с.~5, 15--18]{Wang2006}: лінійна ITG нестійка лише при $0{,}42<r<0{,}76$, але турбулентність \emph{розтікається} (turbulence spreading) у лінійно стійкі зони з порівнянною інтенсивністю. Отже, перенесення нелокальне, і локальне замикання $\chi_{s}(\grad T)$ 1.5D-кодів~--- лише наближення; зональні потоки~--- п.~\ref{sec:06-shear}.
\end{corpusexample}

\section{Домішки}\label{sec:05-imp}

Кожен зарядовий стан $z=0,\dots,Z$ домішки з ядерним зарядом $Z$ описує рівняння неперервності~\cite[с.~4--7, (1), (3)--(5)]{Hulse1992}:
\begin{equation}\label{eq:05-hulse}
  \frac{\pa n_z}{\pa t}=-\frac1r\frac{\pa}{\pa r}(r\Gamma_z)+I_{z-1}n_{z-1}-(I_z+\mathcal{R}_z)n_z+\mathcal{R}_{z+1}n_{z+1}-\frac{n_z}{\tau_z}+S_z,\qquad \Gamma_z=-D_z\frac{\pa n_z}{\pa r}+v_zn_z,
\end{equation}
де $I_z=n_e\fsa{\sigma v}_{\mathrm{ion}}$, а $\mathcal{R}_z$ містить діелектронну й радіаційну рекомбінацію ($\propto n_e$) і перезарядку на нейтральному водні ($\propto n_{\mathrm H^0}$, зокрема на нейтралах пучка). Доданок $n_z/\tau_z$ описує паралельні втрати в SOL. Там $\tau_z=\lambda_{\mathrm{SO}}^2/D$, де $\lambda_{\mathrm{SO}}$~--- довжина спаду в SOL~\cite[с.~10, (11)]{Hulse1992}. Система «радіус $\times$ заряд» блочно-тридіагональна, тож її розв'язують неявно. Якщо $D$, $v$ не залежать від $z$, то в області без джерел нульовий повний потік дає параметризацію пінчу~\cite[с.~9--10, (7)--(9)]{Hulse1992}:
\begin{equation}\label{eq:05-Cv}
  \frac{v}{D}=C_v\frac{\pa\ln n_e}{\pa r}\quad\Longrightarrow\quad n_Z^{\mathrm{tot}}\propto n_e^{C_v}.
\end{equation}
$C_v=1$~--- стала концентрація, $C_v>1$~--- накопичення домішки в центрі. Оскільки $\pa_rn_e<0$, пінч спрямований усередину.

\begin{corpusexample}[германій у TFTR і чутливість до атомних даних]
TFTR, $T_e(0)\approx\SI{6}{кеВ}$; ілюстративні (не виміряні) $D=\SI{e4}{см^2/с}=\SI{1}{м^2/с}$, $C_v=1$, $\lambda_{\mathrm{SO}}=\SI{2}{см}$. Після інжекції Ge за 50--100~мс доходить до центру й іонізується до He-подібного Ge$^{+30}$; значна частина втрачається в SOL~\cite[с.~11, 18]{Hulse1992}. Головний висновок~\cite[с.~22]{Hulse1992}: зміна $D$ удвічі і зміна швидкостей іонізації чи рекомбінації удвічі однаково сильно змінюють еволюцію Ge$^{+25}$. Тож визначити $D$ з таких експериментів можна не точніше, ніж відомі атомні дані.
\end{corpusexample}

\textbf{Неокласика й обертання.} Для важкої домішки (маса $m_I$, заряд $Z_I$) у режимі PS за обертання $\omtor$ Вонг і Ченг отримали (переписано у вигляді повного квадрата за $T_I=T_i$)~\cite[с.~6, (6)]{Wong1987}
\begin{equation}\label{eq:05-wong}
  D_{\mathrm{neo}}=\frac{q^2\nu\,m_IT_I}{(Z_IeB)^2}\Big[1+\frac{m_I\omtor^2R^2}{2T_I}\Big(1-\frac{m_i}{m_I}\frac{Z_IT_e}{T_i+Z_iT_e}\Big)\Big]^2 .
\end{equation}
При $\omtor=0$ це $q^2\rho_I^2\nu$ (Галєєв--Сагдєєв), тобто оцінка PS з п.~\ref{sec:05-coll}. $D_{\mathrm{neo}}$ залежить лише від $\omtor^2$, тож не залежить від напрямку обертання. При $m_I\omtor^2R^2\gg T_I$ він зростає на 1--2 порядки~\cite[с.~6--7]{Wong1987} (обертова система~--- п.~\ref{sec:02-rot}). Неокласичний потік важких домішок на краю спрямований усередину. Його компенсує температурне екранування, і між ELM формується порожнистий профіль W із максимумом на вершині педесталу~\cite[с.~4]{vanVugt2019}.

\textbf{Перенесення W під час ELM.} JOREK з трасуванням повних орбіт W для ELM типу~I в ASDEX Upgrade (\#33616)~\cite[с.~6--11]{vanVugt2019}: W переносить дрейф $\ExB$ у полі пілінг-балонної моди, а магнітне збурення майже нічого не додає. Рух супердифузійний, W перерозподіляється при $\psiN\in[0{,}75;\,1{,}05]$, ефективне $D\approx\SI{11}{м^2/с}$ (з $\fsa{|\grad\psiN|^{-2}}=\SI{1,51}{м^2}$ при $\psiN=0{,}95$). Для добре екранованого (порожнистого) профілю \mbox{10--20\,\%} W проникає \emph{всередину}, тобто ELM можуть підвищувати вміст W у ядрі. 1D-дифузія цієї глибини проникнення не відтворює. Застереження: зіткнень немає, опір $\approx4\times$ спітцерівського.

\section{Скейлінги утримання, режими L і H, педестал}\label{sec:05-scal}

Час утримання $\tauE=W/P$~--- запасена енергія, поділена на потужність нагріву. Стандартна картина (поза корпусом): у режимі H на краю утворюється транспортний бар'єр~--- \emph{педестал} шириною кілька сантиметрів із крутими $\grad p$, які періодично скидають ELM. $\tauE$ зростає в середньому приблизно вдвічі порівняно з L-режимом~\cite[с.~8]{Fukuyama1992}.

\textbf{Безрозмірний скейлінг ELMy H-mode}~\cite[с.~1, (1)--(2)]{Mazzucato2000}:
\[\tauE\Omc{i}\propto\rho_*^{-2{,}70}\beta^{-0{,}90}\nustar^{-0{,}01}M^{0{,}96}\qnf^{-3{,}0}A^{-0{,}73}\kappa^{2{,}3}.\]
Показник $-2{,}70$ лежить між бомівським $-2$ і гіро-бомівським $-3$ \eqref{eq:05-gB}, а від зіткненнєвості залежність практично відсутня.

\textbf{L-скейлінг ITER89 і суперечність корпусу.} TASK/TR нормовано на~\cite[с.~8, (22)]{Fukuyama1992}
\begin{equation}\label{eq:05-iter89}
  \tauE^{\mathrm{ITER89}}=0{,}048\,M^{0{,}5}\Ip^{0{,}85}R^{1{,}2}a^{0{,}3}\kappa^{0{,}5}\bar n^{0{,}1}B^{0{,}2}P^{-0{,}5}\ [\text{с; МА, м, }10^{20}\,\text{м}^{-3},\text{ Тл, МВт}]
\end{equation}
з множником 1,5 в омічній фазі (скромно порівняно з $\approx2$ у режимі H). Модель відтворює залежність від $P$ і (при високій густині) від $\bar n$. Проте \emph{сама модель} дає $\tauE\propto\Btor$ без залежності від $\Ip$, тоді як \eqref{eq:05-iter89} має $\Ip^{0{,}85}B^{0{,}2}$. Автори визнають це межею застосовності моделі дрейфових хвиль~\cite[с.~8]{Fukuyama1992}. Інші джерела корпусу теж не узгоджуються: у~\cite{Mazzucato2000} $\tauE$ залежить від струму сильно, через $\qnf^{-3}\propto(\Ip R/(Ba^2))^3$, а в~\cite{Park1989} слабко, як $\Ip^{0{,}18}$. Тож результати TASK/TR надійні лише за майже сталих $\Ip$ і $\Btor$.

\textbf{Піковість густини (TFTR).} Регресія за 561 точкою ($\Ip=0{,}8$--$1{,}8$~МА, $P_B=10$--30~МВт NBI)~\cite[с.~4--5]{Park1989}: $\tauE=2{,}4\cdot10^{-3}F_{ne}^{0{,}76}P_B^{-0{,}12}\Ip^{0{,}18}$, $F_{ne}=n_e(0)/\fsa{n_e}$. Показники прочитано зі скану (0,76 може бути 0,78), одиниць не вказано, тож префактор непридатний. Розряди з $F_{ne}>2$ не мали пилки й мали кращий $\tauE$.

\textbf{Поріг L--H і педестал у Fenix}~\cite[с.~4]{Fable2022}: перехід відбувається, коли іонний тепловий потік через вершину педесталу досягає $Q_{i,\mathrm{ped}}=0{,}0011\,n^{1{,}07}\Btor^{0{,}76}S$ ($S$~--- площа поверхні). Одиниць не вказано; оцінка для AUG (задача~3) узгоджується з $n$ у $10^{19}$\,м$^{-3}$ і $Q$ у МВт (наш висновок, не джерела). У режимі H тиск на вершині педесталу обмежено: $P_{\mathrm{ped,crit}}=0{,}33R^{-0{,}38}\ee^{5\delta}\Ip^{1{,}25}\kappa^{0{,}62}\betaN^{0{,}43}$. Для всієї потужності~--- $P_{\mathrm{LH}}=3{,}24M^{-1}B^{0{,}75}n_{20}^{0{,}60}R^{0{,}98}a^{0{,}81}$~\cite[с.~3]{Mazzucato2000}.

\begin{practicum}[1D баланс $T_e$, $T_i$ за Тищенком]
\textbf{Задача.} Напишіть неявний розв'язувач \eqref{eq:05-tysh} ($\sim$60 рядків \texttt{numpy}, середовище \texttt{fusion equilibrium challenge/starter/.venv}): циліндр $a=\SI{1,25}{м}$, $R_0=\SI{3}{м}$ (JET, \cite[с.~1]{Fukuyama1992}); $n_e=n_i=n_0[0{,}2+0{,}8(1-x^2)]$, $x=r/a$, $n_0=\SI{3,8e19}{м^{-3}}$; NBI \SI{10}{МВт} з профілем $\propto\ee^{-(x/0{,}4)^2}$, 60\,\% на іони; $P_{ei}$ за \eqref{eq:05-tau}; $P_{\mathrm{br}}=5{,}35\cdot10^{-37}\Zeff n_e^2\sqrt{T_e[\text{кеВ}]}$~Вт/м$^3$, $\Zeff=1{,}5$; $T(a)=\SI{0,1}{кеВ}$; $\chi_{i}=\SI{0,50}{м^2/с}$, $\chi_{e}=\SI{0,27}{м^2/с}$. Інтегруйте до стаціонару. \textbf{Очікуваний результат (наш прогін, сітка 101 і 401 вузол, розбіжність $<0{,}1\,\%$):} $T_{e0}\approx\SI{36,7}{кеВ}$, $T_{i0}\approx\SI{33,2}{кеВ}$, $\tauE\approx\SI{2,0}{с}$. Це в 2,5--3 рази вище, ніж у табл.~4.2. Щоб отримати $T_{i0}\approx\SI{13}{кеВ}$, потрібно помножити обидва $\chi_{s}$ на $f\approx2{,}5$ (60\,\% NBI на іони: $T_{e0}\approx\SI{14,7}{кеВ}>T_{i0}$) або $f\approx2{,}7$ (70\,\% на іони: $T_{e0}\approx\SI{11,9}{кеВ}$).
\textbf{Що дослідити:} (1) знайдіть $f$, за якого $T_{i0}=\SI{13,35}{кеВ}$; (2) з'ясуйте, яка з відмінностей від~\cite{Tyshchenko2021} пояснює множник: витягнутість (об'єм), профіль осадження TRANSP, нестаціонарність розряду, перезарядка; (3) додайте $P_\alpha$ і перевірте тренд «α-нагрів підвищує $T_e$ більше, ніж $T_i$».
\end{practicum}

\subsection*{Контрольні запитання}
\begin{enumerate}
  \item Виведіть \eqref{eq:05-nustar} і покажіть, що $D_{\mathrm b}$, $D_{\mathrm p}$, $D_{\mathrm{PS}}$ зшиваються на межах режимів. Чому бутстреп-струм не потребує зовнішнього поля?
  \item Як в'язкість \eqref{eq:05-mu} поєднує форми Хіршмана і Брагінського? Що втрачається через те, що $\mu_{sj}$~--- функції потоку?
  \item Покажіть, що форми енергетичного потоку з $\tfrac32T\Gamma$ і $\tfrac52T\Gamma$ еквівалентні. Яку роботу тиску при цьому треба перенести в джерело?
  \item Чому 1.5D-модель з локальним $\chi_{s}(\grad T)$ не може описати результат GTC про розтікання турбулентності?
  \item Чому зі спостережень Ge$^{+25}$ не можна надійно визначити $D$~\cite{Hulse1992}? Як ELM можуть \emph{збільшити} вміст W у ядрі~\cite{vanVugt2019}?
\end{enumerate}

\subsection*{Задачі}
\begin{enumerate}
  \item \emph{Режими й аномальність.} Для ядра JET \#41069 ($n_e=\SI{3,8e19}{м^{-3}}$, $T_e=\SI{9,9}{кеВ}$, $T_i=\SI{13,35}{кеВ}$~\cite{Tyshchenko2021}; $q=1$, $R=\SI{3}{м}$, $\varepsilon=0{,}1$, $B=\SI{3,45}{Тл}$~\cite{Fukuyama1992}; D) обчисліть $\lnL$, $\tau_e$, $\tau_i$, $\nu_{*e}$, $\nu_{*i}$, $\rho_i$ і $\chi_{i}^{\mathrm{neo}}$ (прийміть $\nu_{ii}=1/\tau_i$, $v_T=\sqrt{T/m}$). У скільки разів $\chi_{i}=\SI{0,5}{м^2/с}$ Тищенка більший за неокласичний?
  \item \emph{Домішки.} (а) За \eqref{eq:05-Cv}: у скільки разів концентрація в центрі перевищує крайову для $C_v=2$, якщо $n_e(0)/n_e(a)=3$? (б) Знайдіть $\tau_z$ у SOL для прикладу TFTR~\cite{Hulse1992}. (в) За \eqref{eq:05-wong} оцініть підсилення $D_{\mathrm{neo}}$ германію ($Z_I=22$; D, $Z_i=1$, $T_I=T_i$) у TFTR при $r=a/2$: $\omtor=\SI{e5}{рад/с}$, $R=\SI{2,5}{м}$, $T_e=\SI{1,23}{кеВ}$, $T_i=\SI{3,08}{кеВ}$ (див. задачу~3 розд.~\ref{ch:particles}).
  \item \emph{Скейлінги.} (а) За \eqref{eq:05-iter89} обчисліть $\tauE$ для ITER-класу~\cite{Fukuyama1992} ($M=2{,}5$, $\Ip=\SI{22}{МА}$, $R=\SI{6}{м}$, $a=\SI{2}{м}$, $\kappa=2$, $\bar n=1$, $B=\SI{4,85}{Тл}$, $P=\SI{100}{МВт}$) і з множником 1,5. Як змінюється $\tauE$ при подвоєнні $B$ за скейлінгом і за моделлю TASK/TR? (б) Для AUG ($M=2$, $B=\SI{2,5}{Тл}$, $\bar n=\SI{4e19}{м^{-3}}$, $R=\SI{1,65}{м}$, $a=\SI{0,5}{м}$, $\kappa=1{,}7$, $S\approx4\pi^2Ra\sqrt{(1+\kappa^2)/2}$) порівняйте $P_{\mathrm{LH}}$~\cite{Mazzucato2000} з $Q_{i,\mathrm{ped}}$~\cite{Fable2022} для $n$ у $10^{19}$ і у $10^{20}$\,м$^{-3}$. Який вибір одиниць правдоподібний?
\end{enumerate}
